\include{includes}
%SetFonts


\title{Topic 3: Asymptotics}
\author{02-680: Essentials of Mathematics and Statistics}
%\date{}							% Activate to display a given date or no date

\begin{document}
\maketitle

 \section{Big O}
 To understand function behavior of say function $f(x)$ we often want to look at it its growth compared to another function $g(x)$,
 for instance we might ask:
 \begin{tcolorbox}
 \centering
 Does $f(x)$ grow no faster than $g(x)$?
 \end{tcolorbox}
 
 Define the \textbf{set} of functions $f(x)$ for which this is true: \\
 For functions $h: \reals^{\ge 0} \rightarrow \reals^{\ge 0}$ and $g: \reals^{\ge 0} \rightarrow \reals^{\ge 0}$
\[
O\big(g\left(x\right)\big) := \bigg\{h(x) \biggm| \exists x_0 \ge 0, c > 0 : \forall x \ge x_0 : h(x) \le c \cdot g(x) \bigg\}
\]
 (notice the explicit use of $\ge$ vs $>$)

The question above then becomes 
\begin{tcolorbox}
\centering
Is $f(x) \in O\big(g(x)\big)$?
\end{tcolorbox}

\begin{aside}[pedantry detail]
You might see someone use $f(x) = O(g(x))$, this is not technically correct, as its asking if a function is 
equivalent to a set of functions. 
\end{aside}

So as an example: $f(x) := 4x$ and $g(x)=x^2$. 

\begin{center}
\begin{tikzpicture}
\begin{axis}[
     legend style={at={(1,1)},anchor=north west},
    samples=100,smooth,
    ymin=0,ymax=60,
    xmin=0,xmax=8,
    xlabel={$x$},
    ylabel={$f(x)/g(x)$},
    domain=0:8,
]
\addplot [blue] {4*\x}; \addlegendentry{$4x$}
\addplot [red] {\x^2}; \addlegendentry{$x^2$}
\addplot [dotted, no marks] coordinates {(4,0) (4,60)};
\end{axis}
\end{tikzpicture}
\end{center}

\begin{itemize}
\item $f(x) \in O(g(x) + h(x)) \iff f(x) \in O(\max(g(x),h(x)))$
\item $f(x) \in O(g(x))$ and $g(x) \in O(h(x)) \implies f(x) \in O(h(x))$
\item $f(x) \in O(h_1(x))$ and $g(x) \in O(h_2(x)) \implies f(x) + g(x) \in O(h_1(x) + h_2(x))$
\item $f(x) \in O(h_1(x))$ and $g(x) \in O(h_2(x)) \implies f(x) \cdot g(x) \in O(h_1(x) \cdot h_2(x))$
\end{itemize}

\subsection{Polynomials}
Let $p(x) = \displaystyle_{i=0}^k a_ix^k$ be a polynomial of order $k$,
$p(x) \in O(x^k)$.

Basic idea of proof: for $i<k$, $a_ix^i < |a_i|x^k$ for any $n>1$. 

\subsection{Logarithms \& Exponentials}
Log << Polynomial << Exponential

$\log x \in O(n^\epsilon)$ for $\epsilon >0$. (base of the log does not matter)

$p(x) \in O(b^x)$ for any $b>1$ and $k\ge 0$ for the polynomial above.

$b^x \in O(c^x) \iff b \le c$. 

\subsection{Useful Examples}
Is $3x^2+2 \in O(x^2)$? is $x^3 \in O(x^2)$?

Is $3x+4y+5xy \in O(x)$? Is it $\in O(y)$? is it $\in O(xy)$?

\section{Others}
\[
\Omega\big(g\left(x\right)\big) := \bigg\{h(x) \biggm| \exists x_0 \ge 0, d > 0 : \forall x \ge x_0 : h(x) {\color{red}\boldsymbol{\ge}} d \cdot g(x) \bigg\}
\]
 \begin{tcolorbox}
 \[f(x) \in O(g(x)) \iff g(x) \in \Omega(f(x))\]
 \end{tcolorbox}
\[
\Theta\big(g\left(x\right)\big) := \bigg\{h(x) \biggm| h(x) \in O( g(x)) \wedge h(x) \in \Omega(g(x)) \bigg\}
\]
\[
o\big(g\left(x\right)\big) := \bigg\{h(x) \biggm| {\forall \color{red}c>0}  : \exists x_0 \ge 0: \forall x \ge x_0 : h(x) {\color{red}<} c \cdot g(x) \bigg\}  := \bigg\{h(x) \biggm| h(x) \in O( g(x)) \wedge h(x) \notin \Omega(g(x)) \bigg\}
\]
\[
\omega\big(g\left(x\right)\big) := \bigg\{h(x) \biggm| {\forall \color{red}c>0} : \exists x_0 \ge 0 : \forall x \ge x_0 : h(x) {\color{red}>} c \cdot g(x) \bigg\}  := \bigg\{h(x) \biggm| h(x) \notin O( g(x)) \wedge h(x) \in \Omega(g(x)) \bigg\}
\]

\subsection{Useful Example}
What is the relationship between these two functions: 
\[f(x) = x^2 \quad g(x) = \begin{cases}x^3 & x\%2=0\\x & x\%2=1\end{cases}\]
Is $f(x)\in O(g(x))$? Is $f(x)\in \Omega(g(x))$? Is $f(x)\in o(g(x))$? Is $f(x)\in \omega(g(x))$?

%%%%%%%%%%%%%%%%%%
\section*{Useful References}

Liben-Nowell, ``Connecting Discrete Mathematics and Computer Science, 2e''. \S 6.2

\end{document}  